\section{A Coupled Spatial-Economic World-Model Benchmark}\label{sec:setting}

\subsection{Economic Mechanics, Deterministic Replay, and Model Ladder}\label{sec:engine}\label{sec:emulator}

To make discrepancies between imagined and true outcomes exact, we patched the Micropolis C++ core to remove wall-clock seeding, uninitialised pointers, and non-deterministic memory state. Across all \OracleCities{} canonical cities stepped for \OracleMonths{} months with disasters active, full-state hashes match bit-for-bit across fresh processes, dirty simulator instances, and host architectures. Restoring the observation grid, bookkeeping ledger, internal PRNG, and active-sprite state reproduces twelve-month rollouts bit-for-bit. Active vehicle and disaster sprites are omitted from world-model inputs and introduce mild partial observability, mainly by shifting the shared PRNG stream (Appendix~\ref{app:repro}).

Episodes span \Horizon{} monthly steps on a \WorldW{}-by-\WorldH{} tile grid. Each month combines fiscal choices (municipal tax rate up to \TaxMax{} per cent plus transport, fire, and police funding) with up to \MacrosPerMonth{} spatial builds or demolitions (Appendix~\ref{app:data}). Capital spending deducts immediately from reserves, whereas tax receipts and public-service budgets settle on annual cycles. We collect \TrainTransitions{} monthly transitions across \TrainFamilies{} training cities under a \emph{behaviour mix} of scripted builders, stochastic editors, fiscal sweeps, and infrastructure shocks.

Evaluation spans \AllCities{} held-out cities (\AllInstances{} starting instances) from disjoint map seeds. We split these into a \HeldOutFamilies-city core benchmark (\MainInstances{} instances) for compute-intensive scaling, RL, patching, replanning, and coverage experiments, and an \XCities-city expansion split (\XInstances{} instances) for cross-model audits. On each held-out start, executing the behaviour mix yields \PrLoggedPerStart{} open-loop \emph{logged plans} (Appendix~\ref{app:searches}).

Each world model is a stochastic one-month latent-variable ConvNeXt U-Net \citep{ronneberger2015unet,liu2022convnext} with a \ZDim-dimensional variational noise latent. It predicts structural tiles, power and flag planes, and spatial overlays, paired with a bookkeeping network predicting municipal scalars from map census counts (Appendix~\ref{app:emulator}). Our \Ladder-checkpoint ladder crosses three parameter scales (\ParamsFifteen{}M, \ParamsFifty{}M, \ParamsHundred{}M) with two data budgets (\TrainTransitions{} vs.\ stratified \SmallTransitions{} transitions). Each one-step model trains for \HpSamplesOs{} examples (two epochs over the full store, under two H200-equivalent GPU hours per run; Table~\ref{tab:compute}) with mid-training and final snapshots. Each is then fine-tuned via multi-step \emph{pushforward rollout training} \PushforwardCites{} for \HpSamplesSf{} further examples across two snapshots. Our primary model is the full-budget \SizeM{} full-data one-step checkpoint.

\subsection{Short-Horizon Fidelity Gates Mask Task-Horizon Spatial Drift}\label{sec:fidelity}

On a held-out test split (\TestTransitions{} transitions, \TestPlans{} trajectories), all \TestRuns{} budget-matched models\ifTestHalfAllPass{} and all \TestHalfN{} mid-training snapshots (Appendix~\ref{app:ladderfid})\fi{} clear standard short-horizon benchmarks. One-month \emph{structural} changed-tile F1 (quotienting out traffic and animation flicker; Appendix~\ref{app:other}) spans \TestFOneMin{}--\TestFOneMax{} above the \TestBarFOne{} half-oracle bar (Table~\ref{tab:test}). Median \RollA-month relative population and treasury errors stay below \KillPop{}. The primary \SizeM{} checkpoint reaches \TestMainFOne{} \ci{\TestMainFOneCI} structural F1 (engine cross-seed ceiling \TestFloorFOne{}; nearest-neighbour \TestNNFOne{}), with \RollA-month population and funds errors at \TestPopRatio{} and \TestFundsRatio{} times engine seed spread. At this two-epoch budget, held-out variational log-likelihood still improves on all full-data models, whereas all \HlSmallN{} small-data models already overfit relative to mid-training (Table~\ref{tab:horizon}).

Clearing short-horizon gates nevertheless masks substantial drift by month \Horizon{}. Rolled out open loop to month \Horizon{}, the primary model exceeds start-state copying error on six of \HorOverlays{} spatial overlays and ties on a seventh at \HorRestMin{} times copy error (Table~\ref{tab:horizon}). It beats copying only on \HorBestOverlay{} (\HorBestRatio{}) and peaks on pollution at \HorPollEmu{} vs.\ \HorPollEng{} for engine seed spread. Across \HorOpenPlans{} open-loop plans, it miscounts powered zones by \HorPowEmu{} tiles versus \HorPowCopy{} for copying and \HorPowEng{} for engine seed spread\ifHorizonDone{} (\HorLadGeMin{}--\HorLadGeMax{} overlays equal or exceed copy error across all \LadFinals{} final checkpoints; Appendix~\ref{app:ladderfid})\fi. Multi-step pushforward fine-tuning slightly lowers short-horizon F1 in \SfLowerOne{} of \SfPairs{} parent-child pairs and frequently increases planning false acceptance relative to its one-step parent (Tables~\ref{tab:curves} and~\ref{tab:test}; Appendix~\ref{app:ablations}). Even on \BehPlans{} unoptimised logged plans, the primary model falsely rejects \GrowBehFn{} and \SolvBehFn{} of passing \emph{grow} and \emph{solvent} plans while falsely accepting \TreasBehFp{} and \RecovBehFp{} of failing \emph{treasury} and \emph{recover} plans. Extending one-step training sixfold on a single seed (\ExtTrEpochs{} epochs) lifts held-out likelihood yet worsens \Horizon-month spatial drift, as analysed in Section~\ref{sec:extended}.

\subsection{Economic and Infrastructure Planning Tasks and Pipeline Reference}\label{sec:tasks}\label{sec:measure}

We evaluate five \Horizon-month task families pairing each starting city with a calibrated target $\theta$ across five economic regimes (Appendix~\ref{app:tasks}). In \emph{grow} and \emph{solvent}, plans must raise population to at least $\theta$ while keeping liquid reserves non-negative or at least half of initial funds. In \emph{treasury}, second-year annual tax revenue must reach $\theta$ through tax-base growth without deficit spending or rate hikes. In \emph{recover}, a city must absorb a month-one earthquake shock and restore at least $\theta$ of its pre-quake population. In \emph{power}, a city must reach at least $\theta$ developed zones with at least \PowerFrac{} per cent electrified and at most one added power plant. Placing $\theta$ midway between the best constant no-build policy and engine-side search leaves \Instances{} of \InstancesPlanned{} core instances requiring active investment (\WellPosed{} \emph{well-posed} with at least \WellPosedSE{} standard errors of headroom; Table~\ref{tab:tasks}).

Single-seed engine optimisation passes its own seed on \NineSeedOwnLo{}--\NineSeedOwnHi{} per cent of instances yet transfers to fresh seeds on only \NineSeedOtherLo{}--\NineSeedOtherHi{} per cent. Ground truth is therefore majority passage across \OracleSeeds{} reserved oracle seeds (at least sixteen of \OracleSeeds{} passing). Search candidates are ranked by mean continuous \emph{margin} over \KSearch{} rollouts and accepted by majority over \KVerdict{} fresh world-model rollouts (at least four of \KVerdict{} passing). We track three core quantities.

\begin{itemize}
\item \textbf{False acceptance.} The share of instances where the judge accepts the selected plan yet the \OracleSeeds-seed oracle rejects it.
\item \textbf{Ground-truth pass rate.} The share of instances whose selected plan passes the \OracleSeeds-seed oracle.
\item \textbf{False rejection.} The share of truly passing selected plans that the judge rejects.
\end{itemize}

Running the identical optimiser with matched rollout budgets ($\KSearch{}$, $\KVerdict{}$) inside an \emph{exact-simulator reference} on fresh non-oracle C++ seeds establishes our baseline for stochastic selection luck. Instance-paired differences define \emph{excess false acceptance} and \emph{ground-truth success shortfall}, capturing the combined operational pipeline effect of model bias, rollout variance, and altered plan rankings (Appendix~\ref{app:stats}; Tables~\ref{tab:main} and~\ref{tab:family}).
